% ======================================================================
%  Spatial Digital Twin for Unsupervised Wildfire Detection
%  Using GOES-16 ABI: Multi-Band Pipeline, Temporal Assimilation,
%  Multi-View Consensus, and Agentic Orchestration
% v3.0 — Expanded for A1-level journal submission
% ======================================================================

\documentclass[11pt,twocolumn,a4paper]{article}

% ---- PdfLaTeX preamble (no fontspec needed) ----
\usepackage[T1]{fontenc}
\usepackage[utf8]{inputenc}

% ---- Language ----
\usepackage[portuguese,english]{babel}
\usepackage[hyphens]{url}

% ---- Math ----
\usepackage{amsmath,amssymb,amsfonts}
\usepackage{bm}

% ---- Graphics & Color ----
\usepackage{graphicx}
\usepackage[table]{xcolor}

% ---- Tables ----
\usepackage{booktabs}
\usepackage{tabularx}
%\usepackage{multirow}  % not available in basic TeX Live

% ---- Bibliography ----
\usepackage[numbers,sort&compress]{natbib}

% ---- Algorithms ----
%\usepackage{algorithm}    % not available in basic TeX Live
%\usepackage{algpseudocode} % not available in basic TeX Live

% ---- Links ---- 
\usepackage[colorlinks=true,linkcolor=blue,citecolor=blue,urlcolor=blue]{hyperref}

% ---- Page layout ----
\usepackage[top=2.5cm,bottom=2.5cm,left=2.0cm,right=2.0cm]{geometry}
\setlength{\columnsep}{0.8cm}

% ---- Custom commands ----
\newcommand{\method}[1]{\texttt{#1}}
\newcommand{\code}[1]{\texttt{#1}}
\newcommand{\abbrev}[1]{{\small\textsc{#1}}}
\newcommand{\N}{\mathbb{N}}
\newcommand{\R}{\mathbb{R}}
\newcommand{\C}{\mathbb{C}}
\newcommand{\E}{\mathbb{E}}
\newcommand{\sK}{\mathcal{K}}
\newcommand{\sL}{\mathcal{L}}
\newcommand{\sG}{\mathcal{G}}
\newcommand{\prob}{\text{prob\_or}}
\newcommand{\ru}{\text{ru}}
\newcommand{\clip}{\text{clip}}

\hyphenation{di-gi-tal twin fire-line Koop-man Pi-GNN}

% ======================================================================
\title{\Large\textbf{Spatial Digital Twin for Unsupervised Wildfire Detection\\[4pt]
with GOES-16 ABI: Multi-Band Pipeline, Temporal Assimilation,\\[4pt]
Multi-View Consensus, and Agentic Orchestration}}

\author{Nauber Gois$^{1,\dagger}$ \\
\small $^{1}$Laboratory of Digital Twins / qclawmonitor Project, Cear\'{a}, Brazil \\
\small $\dagger$ Corresponding author: \texttt{nauber.gois@qclawmonitor.dev} \\
\small \url{https://github.com/naubergois/ceara-queimadas}}

\date{June 7, 2026\vspace{4pt}\\
\small v3.1 --- Updated literature survey (GraphFire-X, Thermodynamic GeoAI,
FireCastNet), FIRMS validation, expanded discussion}

% ======================================================================
\begin{document}

\maketitle

% ======================================================================
\begin{abstract}

and autonomous agentic decision support.
Code available at \url{https://github.com/naubergois/ceara-queimadas}.
\end{abstract}

%% Portuguese abstract
\begin{otherlanguage}{portuguese}
\begin{abstract}
A detecção de incêndios florestais nos biomas brasileiros---Amazônia, Cerrado, Pantanal e Caatinga---continua sendo um desafio crítico, sem gêmeos digitais operacionais que integrem dados de satélite em tempo real, simulação física e suporte autônomo à decisão. Este artigo apresenta um framework não supervisionado de gêmeo digital espacial para detecção de queimadas usando dados GOES-16 ABI L2 CMIPF sobre o estado do Ceará, Brasil. O sistema combina: (i) filtragem mediana multi-escala da banda térmica T\_B13 (10,3~$\mu$m) com contraste espectral ponderado T\_B7$-$T\_B14 (3,9$-$11,2~$\mu$m); (ii) assimilação temporal via fusão probabilística com decaimento exponencial (prob\_or); (iii) um método de persistência combinada fundindo intensidade de pico, média temporal e razão de ativação entre granules horários; (iv) um ensemble de consenso exigindo concordância de 2 em 3 visões físicas (gêmeo digital, persistência, residual espacial); e (v) um pipeline de orquestração agentiva baseado em LangGraph com raciocínio ReAct, cruzamento geoespacial e geração auditável de alertas. A avaliação contra 76 focos de referência do INPE em 31/10/2024 com calibração automática de contaminação maximizando F1 produz uma hierarquia de desempenho consistente: gêmeo digital (P=0,032, R=0,106, F1=0,049), isolation forest (P=0,026, R=0,063, F1=0,037), persistência combinada (P=0,031, R=0,007, F1=0,012), residual espacial (P=0,019, R=0,004, F1=0,006) e ensemble de consenso (P=0,000, F1=0,000). Benchmarks sintéticos confirmam que todos os métodos excedem precisão $>$ 0,8 e F1 $>$ 0,8 sob condições alinhadas (8/8 testes aprovados). Identificamos quatro desalinhamentos estruturais---temporal, semântico, de escala e de produto---como causas raiz da lacuna sintético-real. Estendendo esta base diagnóstica, introduzimos o framework matemático Neural Koopman Physics-Informed Graph Neural Network (NeKo-PIGNN) para modelagem preditiva de propagação de fogo, combinando a teoria do operador de Koopman (lacuna absoluta na literatura de incêndios florestais) com GNNs com restrição física regularizadas pela equação de Rothermel. Um pipeline LangGraph orquestra 9 agentes de IA especializados para coleta de dados, análise geoespacial, avaliação de risco climático, análise de persistência GOES-16, fusão de evidências, diagnóstico ReAct, geração de alertas, relatórios técnicos e auditoria---tudo integrado com um frontend React/TypeScript e banco de dados espacial PostGIS. Propomos um roteiro completo de seis camadas para um gêmeo digital bidirecional de incêndios brasileiro integrando INPE, NASA FIRMS, MapBiomas, redes de sensores IoT e suporte autônomo à decisão baseado em agentes.
Código disponível em \url{https://github.com/naubergois/ceara-queimadas}.\end{abstract}
\end{otherlanguage}

\noindent\textbf{Keywords:} Digital Twin; Wildfire Detection; GOES-16; Unsupervised Learning; Remote Sensing; Brazilian Biomes; Multi-View Consensus; Koopman Operator; Physics-Informed GNN; LangGraph; Agentic AI \\
\textbf{Palavras-chave:} Gêmeo Digital; Detecção de Queimadas; GOES-16; Aprendizado Não Supervisionado; Sensoriamento Remoto; Biomas Brasileiros; Consenso Multi-Vista; Operador de Koopman; GNN com Informação Física; LangGraph; IA Agêntica

% ======================================================================
\section{Introduction}
\label{sec:intro}

\subsection{Context and Motivation}

Wildfires in Brazil pose a growing threat to globally significant ecosystems. The Amazon---Earth's largest terrestrial carbon sink---recorded over 110,000 fire foci in 2024 alone~\cite{inpe2025}. The Cerrado, the most biodiverse savanna on the planet, concentrates the highest number of active fire detections nationally, with recurrent fires destroying vast tracts of native vegetation annually. The Pantanal suffered catastrophic fires in 2020 and 2024, destroying over 30\% of the biome~\cite{mapbiomas2025}. In the Caatinga, an exclusively Brazilian semi-arid biome covering approximately 70\% of the state of Cear\'{a}, anthropogenic fires for land management and agricultural expansion are increasingly severe during prolonged drought periods.

Despite this urgency, Brazil does not currently operate a digital twin for wildfire detection and response~\cite{morsali2026,raha2025}. While the concept of digital twins was originally formulated for manufacturing~\cite{grieves2014,tao2019,barricelli2019}, recent advances have extended it to environmental systems~\cite{rasheed2020,fuller2020,jones2020}. International Earth system digital twin initiatives~\cite{bauer2021_nature,voosen2020,nativi2021,bauer2023_natrev} demonstrate the feasibility of bidirectional virtual replicas for environmental monitoring, yet these have not been deployed for operational wildfire management in Brazil. Existing monitoring systems---\abbrev{INPE} Queimadas, \abbrev{PRODES}, \abbrev{DETER}, \abbrev{BDQueimadas}, and MapBiomas Fogo---provide essential fire alert and deforestation monitoring services. However, these are fundamentally read-only monitoring platforms, not bidirectional digital twins as defined by \abbrev{ISO}/\abbrev{IEC} 23247~\cite{iso2021}. A systematic survey of digital twin technologies for fire detection in Brazilian biomes~\cite{mukkavilli2023,prapas2023} confirms that no operational digital twin exists for Brazilian wildfire detection---all current programs are alert/monitoring systems lacking predictive simulation, what-if scenario modeling, bidirectional control loops, and \abbrev{AI}-driven autonomous decision support. The fire activity in the Cerrado and Amazon biomes has been extensively documented by satellite-based studies~\cite{barbosa2024_amazon,cano2024_latam}, which further highlight the need for predictive capabilities beyond basic monitoring.

The state of Cear\'{a}, in northeastern Brazil, presents a compelling case study. Its territory spans 148,920~km$^2$ with diverse eco-regions from the humid coastal strip (Mata Atl\^{a}ntica remnants) to the semi-arid Sert\~{a}o (Caatinga biome) and the chapadas of the Ibiapaba and Araripe plateaus. Fire activity peaks between August and December, coinciding with the dry season when relative humidity drops below 30\% and cumulative rainfall deficits exceed 200~mm. The availability of near-real-time \abbrev{GOES-16} data, operational \abbrev{INPE} fire foci records, and a well-defined administrative geography makes Cear\'{a} an ideal proving ground for a spatial digital twin prototype.

\subsection{Related Work}

International research has advanced digital twin-based wildfire management through several distinct architectures:

\begin{itemize}
  \item \textbf{\abbrev{IVSR}}~\cite{morsali2026}: A bidirectional digital twin with autonomous \abbrev{AI} agents, integrating multi-sensor imagery, meteorological data, and 3D forest models with similarity-based scenario alignment. Demonstrates significant reductions in detection-to-intervention latency through closed-loop \abbrev{UAV} and team reallocation.
  \item \textbf{\abbrev{FIRE-VLM}}~\cite{webb2026}: The first vision-language-model-guided reinforcement learning framework in a physics-grounded fire digital twin, achieving up to $6\times$ detection speedup using \abbrev{UAV}s with dual-view sensing. The \abbrev{GIS}-to-simulation pipeline enables automatic construction of fire digital twins from remote sensing data.
  \item \textbf{\abbrev{FIRETWIN}}~\cite{raha2025}: A cyber-physical digital twin using Unreal Engine + \abbrev{CAWFE} model for King Fire reconstruction, providing immersive interactive analysis and multi-modal sensing integration.
  \item \textbf{\abbrev{AIMNET}}~\cite{zhou2025}: An IoT-empowered digital twin for continuous gas emission monitoring with coupled weather-transport modeling, using a three-layer architecture (physical world $\to$ bidirectional links $\to$ digital world).
  \item \textbf{\abbrev{H-RDT}}~\cite{shen2025}: A hazard-responsive digital twin using Physics-Informed Neural Networks for the fire-power-outage-urban-heat-wave cascade.
  \item \textbf{\abbrev{PINN}s for fire spread}~\cite{vogiatzoglou2024,mao2024_pinn}: Physics-informed neural networks learning wildfire propagation parameters~\cite{vogiatzoglou2024} and physics-constrained wildfire spread simulation~\cite{mao2024_pinn}, providing complementary foundations for data-driven digital twins with physical consistency. The broader PINN literature establishes foundational techniques for embedding physical laws into neural architectures~\cite{karniadakis2021,cuomo2022,lu2021_deepxde,markidis2021}, while addressing training pathologies such as gradient stiffness~\cite{wang2021_gradient,krishnapriyan2021} and providing practical training guidelines~\cite{wang2023_expert}.\\
  \item \textbf{Neural ODEs}~\cite{chen2018}: Continuous-time dynamics modeling that naturally accommodates irregularly timed satellite observations. Related reduced-order modeling approaches such as SINDy~\cite{brunton2016_sindy} provide interpretable equation discovery from sparse data.\\
  \item \textbf{Koopman operator theory}~\cite{koopman1931,brunton2021}: A framework for linearizing nonlinear dynamics through observable functions, enabling spectral analysis and optimal control. Key advances include the spectral analysis of fluid flows~\cite{rowley2009,mezic2013}, Dynamic Mode Decomposition with control~\cite{proctor2016}, deep learning for universal linear embeddings~\cite{lusch2018,takeishi2017}, physics-informed DMD variants~\cite{baddoo2023}, and deep dynamical modeling of unsteady flows~\cite{morton2019}. None of these have been applied to wildfire dynamics prior to our work.\\
  \item \textbf{GraphFire-X}~\cite{esparza2025}: A physics-informed graph attention network ensemble for building-scale wildfire vulnerability at the wildland-urban interface (WUI). Combines dual-expert GNN (fire physics-informed convection, radiation, ember probability attention) with XGBoost for structural fragility assessment. Validated on the 2025 Eaton Fire. The physics-informed GNN approach is directly comparable to our NeKo-PIGNN framework, though GraphFire-X targets building vulnerability rather than satellite-based fire propagation.
  \item \textbf{Thermodynamic GeoAI}~\cite{lim2026}: A thermodynamic-inspired explainable GeoAI framework using statistical physics equilibrium (Burden/Capacity) for detecting wildfire-induced PM2.5 anomalies. Validated on the 2023 Canadian wildfires. Provides a complementary approach for smoke plume analysis in Brazilian biomes, identifying phase transitions in spatial systems and causal regime shifts.
  \item \textbf{FireCastNet}~\cite{michail2025}: An Earth-as-a-Graph framework combining 3D Conv with GraphCast GNN for global seasonal fire prediction. Validated on the SeasFire dataset---a multivariate Earth System data cube. Outperforms GRU, Conv-LSTM, U-TAE, and TeleViT baselines, with strong results over South America. It is the first GraphCast-based fire model validated globally and provides the only operational-scale seasonal fire prediction framework applicable to Brazilian biomes.
  \item \textbf{OSMGraphCLIP}~\cite{michail2026_osm}: A CLIP-style geospatial representation model learning global location embeddings from OpenStreetMap graph data. Provides complementary geo-contextual features for fire risk assessment at municipal resolution, scalable to the 184 municipalities of Cear\'{a}.
  \item \textbf{Data-driven Patagonia fire spread}~\cite{becerra2026}: A reaction-diffusion-convection model for forest fire spread in Patagonia, Argentina, using data-driven approaches with over 50,000~ha burned in 2025. Demonstrates the applicability of physics-informed data-driven fire models across South American ecosystems.
\end{itemize}

None of these systems have been adapted to Brazilian biomes or integrated with operational \abbrev{INPE} data flows. The specific challenges of Brazilian biomes---Amazon cloud cover, continental scale (8.5M~km$^2$), limited connectivity in remote areas, and the need to distinguish controlled burns from criminal fires---remain unaddressed by existing international digital twin frameworks. Recent work specific to Brazilian biomes includes deep learning fire detection on Landsat/Sentinel-2 imagery over the Amazon~\cite{sun2024}, network science analysis of Brazilian wildfire patterns~\cite{marli2024}, multi-sensor burned area assessment in the Pantanal~\cite{pereira2024_fire}, and high-resolution digital twin wildfire simulation using 1~m \abbrev{DEM} data~\cite{lee2026}, yet none proposes a complete digital twin architecture with bidirectional control. Broader surveys of deep learning for wildfire detection and risk assessment~\cite{ba2022_wildfire,de2023_cnn} further confirm that hybrid physics-ML approaches---such as our NeKo-PIGNN---remain unexplored. Satellite-based fire trend analysis at the global scale~\cite{andela2023_human} and burned area mapping for climate assessment~\cite{lizundia2024_fire} provide the larger context for our Brazilian case study.

\subsection{Contributions}

The contributions of this work span four domains---detection methodology, mathematical modeling, agentic orchestration, and operational deployment:

\begin{enumerate}
  \item \textbf{Unspervised multi-scale thermal anomaly detection} combining median filtering residuals at three scales (5$\times$5, 9$\times$9, 15$\times$15 pixels) in \abbrev{T\_B13} with weighted spectral contrast (\abbrev{T\_B7}$-$\abbrev{T\_B14}), enabling fire detection without labeled training data.
  \item \textbf{Temporal assimilation framework} (\code{prob\_or}) fusing hourly risk scores with exponential decay, modeling thermal anomaly persistence between \abbrev{GOES-16} scans.
  \item \textbf{Combined persistence method} fusing peak intensity, temporal mean, and activation dwell ratio with adaptive per-hour percentiles and conditional morphological filtering.
  \item \textbf{Consensus ensemble} requiring 2-of-3 agreement among digital twin, persistence, and spatial residual physical views.
  \item \textbf{\abbrev{NeKo-PIGNN} mathematical framework} introducing Neural Koopman Physics-Informed Graph Neural Networks for predictive fire propagation---a combination of Koopman operator theory (absolute literature gap for wildfire) with physics-constrained \abbrev{GNN}s regularized by the Rothermel equation.
  \item \textbf{LangGraph agentic orchestration pipeline} with 9 specialized \abbrev{AI} agents for data collection, validation, geospatial analysis, climate risk, \abbrev{GOES-16} persistence, evidence fusion, ReAct diagnosis, alert generation, and auditing.
  \item \textbf{Systematic diagnosis} of four structural misalignments between \abbrev{GOES-16} \abbrev{CMIPF} grids and \abbrev{INPE} point foci.
  \item \textbf{Six-layer roadmap} for a bidirectional Brazilian fire digital twin integrating satellite data, physical simulation, IoT sensors, 3D visualization, and autonomous agentic decision support.
\end{enumerate}

% ======================================================================
\section{Methodology}
\label{sec:method}

\subsection{Data Sources and Preprocessing}

We use the \abbrev{GOES-16} \abbrev{ABI} L2 \abbrev{CMIPF} (Cloud and Moisture Imagery Product) in netCDF format from the \abbrev{NOAA} AWS Open Data bucket~\cite{noaa2025}. The \abbrev{CMIPF} product provides calibrated brightness temperatures at 2~km resolution at nadir. For fire detection, the multi-spectral capability of \abbrev{GOES-16} \abbrev{ABI} has been extensively validated against higher-resolution sensors such as \abbrev{VIIRS} and \abbrev{MODIS}~\cite{zhang2022_goes16}, and dedicated \abbrev{FRP} retrieval algorithms have been developed for the GOES-R series~\cite{xu2020_goesfrp}. The foundational \abbrev{FRP} methodology derives from early work with experimental satellites~\cite{wooster2003_frp}, while operational active fire products from \abbrev{MODIS}~\cite{giglio2016_modis} and \abbrev{VIIRS}~\cite{schroeder2014_viiirs} provide the de facto standards for validation. Comparative studies of geostationary fire detection algorithms~\cite{hally2017_fire} inform our approach to multi-band thermal anomaly scoring. For the Cear\'{a} region (bbox: min\_lat $= -7.9^\circ$, max\_lat $= -2.5^\circ$, min\_lon $= -41.5^\circ$, max\_lon $= -37.0^\circ$), this yields approximately $72\times 72$ grid cells per channel after remapping.

Three spectral channels are employed:

\begin{table}[htbp]
\centering
\caption{\abbrev{GOES-16} \abbrev{ABI} spectral channels used.}
\label{tab:channels}
\small
\resizebox{\columnwidth}{!}{%
\begin{tabular}{ccl}
\toprule
\textbf{Band} & \textbf{Wavelength} & \textbf{Purpose} \\
\midrule
Band 7  & 3.9~$\mu$m (SWIR)     & Sensitive to sub-pixel hot spots \\
Band 13 & 10.3~$\mu$m (clean IR) & Primary temperature channel     \\
Band 14 & 11.2~$\mu$m (longwave) & Atmospheric window channel      \\
\bottomrule
\end{tabular}
}
\end{table}

The \abbrev{DQF} (Data Quality Flag) field retains only pixels with \abbrev{DQF}=0 (nominal quality). For UTC hours $h \in \{16, 17, 18\}$ on evaluation date 2024-10-31, we download the nearest available scan granule per channel.

Reference fire foci are obtained from the \abbrev{INPE} \abbrev{BDQueimadas} \abbrev{API}~\cite{inpe_bdqueimadas2025} for Cear\'{a} (UF code 23). The \abbrev{INPE} set contains 76 fire foci on 2024-10-31, each with latitude, longitude, and UTC timestamp. After binning to the $72\times 72$ \abbrev{GOES} grid, these map to 49 raw truth cells and 284 truth cells after one iteration of $3\times 3$ morphological dilation. This dilation is a necessary compensation for the coarse \abbrev{GOES} grid resolution ($\sim$0.075$^\circ$ per cell, $\approx$56~km$^2$ per cell) relative to the nearly point-like \abbrev{INPE} foci.

Additionally, the system ingests real-time data from \abbrev{NASA} \abbrev{FIRMS} (\abbrev{VIIRS} and \abbrev{MODIS})~\cite{schroeder2014_viiirs,giglio2016_modis}, Open-Meteo for meteorological variables, and Nominatim for reverse geocoding of fire foci to municipal boundaries. On 2026-06-06, for instance, \abbrev{FIRMS} recorded 28 active hotspots across Cear\'{a} (VIIRS SNPP=4, NOAA20=8, NOAA21=11, MODIS=5) with a mean \abbrev{FRP} of 4.8~MW, while the preceding day (2026-06-05) recorded 142 hotspots (VIIRS SNPP=39, NOAA20=53, NOAA21=44, MODIS=6) with mean \abbrev{FRP}=11.2~MW and 24 severe events (\abbrev{FRP} $\ge$ 20~MW). This day-to-day variability---a 426\% increase driven largely by the Oeiras/Picos cluster (68 foci, 1223.7~MW total \abbrev{FRP})---highlights the dynamic nature of fire activity in the region and validates the need for continuous autonomous monitoring, consistent with the diurnal cycle characterization enabled by geostationary \abbrev{FRP} retrievals~\cite{xu2020_goesfrp,zhang2022_goes16}. The system's capacity for cross-sensor validation across four independent satellite platforms is demonstrated.

\subsection{Multi-Scale Anomaly Score}

For each hourly granule, we compute a continuous anomaly score $s_t(x,y) \in [0,1]$ for each valid cell $(x,y)$ at time $t$:

\begin{equation}
s_t = \tfrac13 \sum_{k\in\{5,9,15\}} \!
  \ru\bigl(\max(0, T_{B13} - M_k(T_{B13}))\bigr)
\label{eq:multiscale}
\end{equation}

where $\ru(x)$ denotes robust normalization to $[0,1]$, $M_k$ is a $k \times k$ median filter producing a local background estimate, and:

\begin{equation}
\ru(x) = \clip\Bigl(\frac{x - P_3(x)}{P_{97}(x) - P_3(x) + \varepsilon}, 0, 1\Bigr)
\label{eq:robust_unit}
\end{equation}

The multi-scale approach captures fires of varying spatial extent: small (\(5\times5\), \(\sim\)10~km diameter), medium (\(9\times9\), \(\sim\)18~km), and large fire fronts (\(15\times15\), \(\sim\)30~km). When both bands 7 and 14 are available, the spectral contrast $\Delta T = T_{B7} - T_{B14}$ is computed with baseline removal at the 55th percentile and fused with weight $w_{\text{BTD}} = 0.55$. The spectral contrast exploits the Planck function's sensitivity to sub-pixel hot spots: fire-affected pixels exhibit elevated brightness temperature at 3.9~$\mu$m relative to 11.2~$\mu$m, yielding positive $\Delta T$ values even when the fire occupies only a fraction of the sensor's instantaneous field of view.

\subsection{Digital Twin Temporal Assimilation}

The spatial digital twin maintains a risk field $R_t \in [0,1]^{H \times W}$ evolving according to the \prob\ fusion rule:

\begin{equation}
R_{t+1} = 1 - (1 - \rho \cdot R_t) \odot (1 - S_{t+1})
\label{eq:probor}
\end{equation}

where $\rho$ is the persistence decay factor (default $\rho = 0.5$), $S_{t+1}$ is the instantaneous anomaly score at time $t+1$, and $\odot$ denotes element-wise multiplication. This probabilistic formulation models the intuition that fire risk persists between \abbrev{GOES-16} scans ($\sim$10~minute granule intervals) but decays as fires may have moved, been extinguished, or obscured by clouds.

The binary fire mask is derived from $R_t$ via a global percentile: cells exceeding the $(1 - \text{contamination})$ percentile of $R_t$ are marked as fire.

\subsection{Combined Persistence Method}

The combined persistence method fuses three temporal features over $T$ hourly snapshots:

\begin{itemize}
  \item \textbf{Peak intensity} ($\hat{S}_{\max}$): $\ru(\max_t s_t)$
  \item \textbf{Temporal mean} ($\hat{S}_{\text{mean}}$): $\ru(\frac{1}{T} \sum_t s_t)$
  \item \textbf{Persistence fraction} ($p$): $\frac{1}{T} \sum_t A_t$ (where $A_t$ is the activation mask at hour $t$)
\end{itemize}

The combined score $C$ is:

\begin{equation}
C = w_p \cdot \hat{S}_{\max} + w_m \cdot \hat{S}_{\text{mean}} + w_r \cdot p
\label{eq:combined}
\end{equation}

with default weights $w_p = 0.42$, $w_m = 0.21$, $w_r = 0.37$ (optimized for precision-recall trade-off through grid search on synthetic data). An adaptive per-hour percentile (base $P_{88}$, adjustable between $P_{79}$ and $P_{93}$ based on score spread) marks cells as ``active'' per hour. When $T \geq 2$, the persistence gate requires $p \geq \max(0.16, 0.56/T)$. Spatial components with fewer than 3 cells are removed via binary\_opening morphology.

\subsection{Consensus Ensemble (Multi-View)}

The consensus ensemble combines three independent physical views to reduce the false positive rate:

\begin{table}[htbp]
\centering
\caption{Three views of the consensus ensemble.}
\label{tab:views}
\small
\resizebox{\columnwidth}{!}{%
\begin{tabular}{cl}
\toprule
\textbf{View} & \textbf{What it measures} \\
\midrule
Digital twin      & Assimilated state across hours (\code{prob\_or}, smoothed risk) \\
Combined persistence & Peak + temporal mean + activation persistence \\
Spatial residual  & \abbrev{T\_B13} anomaly vs.\ local background + \abbrev{BTD} (7$-$14) \\
\bottomrule
\end{tabular}
}
\end{table}

The final mask is obtained by majority voting:

\begin{equation}
M_{\text{ensemble}} = \mathbf{1}[\,(M_{\text{twin}} + M_{\text{persist}} + M_{\text{resid}}) \geq K\,]
\label{eq:consensus}
\end{equation}

with $K = 2$ (2-of-3 agreement) by default. This is strictly unsupervised at inference time and acts as a natural false-positive filter: only pixels that are anomalously hot according to multiple independent physical rationales are flagged.

\subsection{Evaluation Protocol}

Evaluation is performed on a $72\times 72$ grid covering Cear\'{a}. Ground truth comprises \abbrev{INPE} foci aggregated via 2D binning, with optional morphological dilation ($\texttt{--truth-dilate} \in \{0,1,2\}$). Metrics are:

\begin{itemize}
  \item Precision $= \text{TP} / (\text{TP} + \text{FP})$
  \item Recall $= \text{TP} / (\text{TP} + \text{FN})$
  \item F1 $= 2 \cdot P \cdot R / (P + R)$
  \item IoU $= \text{TP} / (\text{TP} + \text{FP} + \text{FN})$
\end{itemize}

Automatic contamination calibration searches for the value maximizing F1 over 32 uniform steps in $[0.004, 0.18]$. \textbf{Methodological note:} this calibration uses \abbrev{INPE} data from the same day for benchmarking and must not be confused with pure unsupervised operation.

% ======================================================================
\section{NeKo-PIGNN: Neural Koopman Physics-Informed Graph Neural Networks}
\label{sec:neko}
% ======================================================================

The four structural misalignments identified in our detection pipeline---temporal, semantic, scale, and product---motivate a fundamentally different approach to fire propagation modeling. We introduce the Neural Koopman Physics-Informed Graph Neural Network (\abbrev{NeKo-PIGNN}) framework, which combines three mathematical innovations that are, individually and jointly, absent from the wildfire literature.

\subsection{Koopman Operator Theory for Fire Dynamics}

The Koopman operator~\cite{koopman1931} provides a representation of nonlinear dynamical systems through an infinite-dimensional linear operator acting on observable functions, with broad applications from fluid dynamics~\cite{mezic2013,rowley2009} to climate variability~\cite{lorenzo2026}. Let $F: M \to M$ be the (unknown) nonlinear map governing fire state evolution on a compact state space $M \subset \R^d$. For any observable function $g: M \to \C$, the Koopman operator $\sK$ is defined as:

\begin{equation}
(\sK g)(z) = g(F(z))
\label{eq:koopman_def}
\end{equation}

The key insight is that while $F$ is nonlinear on $M$, $\sK$ is linear on the space of observables, enabling spectral analysis of the fire dynamics. Evolving the system forward $n$ steps corresponds to:

\begin{equation}
g(z_{t+n}) = \sK^n g(z_t)
\label{eq:koopman_iter}
\end{equation}

For fire propagation, the state vector $z \in M$ includes: cell-wise brightness temperature $T_{B13}$, fire radiative power (\abbrev{FRP}), fuel moisture content, wind vector components, relative humidity, and historical fire recurrence frequency. The observable functions $g$ are learned via a variational autoencoder that projects the high-dimensional state onto a latent space of Koopman eigenfunctions.

\subsection{Extended Dynamic Mode Decomposition with Learned Observables}

Practical approximation of $\sK$ uses Extended Dynamic Mode Decomposition (\abbrev{EDMD})~\cite{williams2015} and its variants~\cite{proctor2016,baddoo2023}. Given snapshots $X = [z_1, \ldots, z_T]$ and $Y = [z_2, \ldots, z_{T+1}]$, and a dictionary of observables $\Theta(z) = [\theta_1(z), \ldots, \theta_d(z)]^\top$, we solve:

\begin{equation}
\sK \approx G^{+} A, \quad G = \Theta(X)^\top \Theta(X), \quad A = \Theta(X)^\top \Theta(Y)
\label{eq:edmd}
\end{equation}

where $G^{+}$ denotes the Moore-Penrose pseudoinverse. The eigenvectors (Koopman eigenfunctions) and eigenvalues (Koopman spectrum) reveal the intrinsic modes of fire propagation---growth rates, oscillation frequencies, and spatial coherence patterns.

Crucially, we learn $\Theta$ directly from satellite data using a variational autoencoder whose latent representation is constrained to approximate Koopman eigenfunctions~\cite{lusch2018,takeishi2017}. The training objective is:

\begin{equation}
\min_{\Theta, \sK} \| \Theta(Y) - \sK \Theta(X) \|_F^2 + \lambda \| \Theta(X) - \hat{\Theta}(X) \|_2^2
\label{eq:vae_koopman}
\end{equation}

where the first term enforces linear dynamics in the latent space and the second ensures the observables reconstruct the input state. This addresses the ``absolute gap'' identified in our research survey: no published work applies Koopman theory to wildfire dynamics~\cite{brunton2021}.

\subsection{Physics-Informed Graph Neural Network Regularization}

The Koopman linearization alone does not guarantee physical consistency---the learned dynamics $z_{t+1} = \Theta^{-1}(\sK \Theta(z_t))$ may violate thermodynamic constraints. We embed physical knowledge via a \abbrev{GNN} operating on the irregular municipal adjacency graph of Cear\'{a}'s 184 municipalities.

Let $\sG = (V, E)$ be a graph where nodes $V$ represent GOES-16 grid cells (or municipal aggregates) and edges $E$ encode spatial adjacency plus wind-direction-weighted propagation corridors. Each node $v$ carries features $h_v^{(0)} = [T_{B13}, \text{FRP}, \text{NDVI}, \text{wind}_x, \text{wind}_y, \text{humidity}, \text{slope}]$. The \abbrev{GNN} message-passing layer is:

\begin{equation}
h_v^{(l+1)} = \sigma\Bigl(W^{(l)} h_v^{(l)} + \sum_{u \in \mathcal{N}(v)} \phi^{(l)}(h_u^{(l)}, e_{uv})\Bigr)
\label{eq:gnn_message}
\end{equation}

where $\mathcal{N}(v)$ are neighbors of $v$, $e_{uv}$ are edge features (distance, wind exposure), and $\phi$ is a learned message function.

The physics-informed loss incorporates the Rothermel fire spread equation~\cite{rothermel1983} as a regularizer, following the Physics-Informed Machine Learning paradigm~\cite{karniadakis2021,cuomo2022}:

\begin{equation}
\mathcal{L}_{\text{total}} = \mathcal{L}_{\text{data}} + \lambda_1 \mathcal{L}_{\text{PDE}} + \lambda_2 \mathcal{L}_{\text{IC}}
\label{eq:pinn_loss}
\end{equation}

This residual directly encodes the reaction-diffusion nature of fire spread, with known training challenges in PINNs~\cite{wang2021_gradient,krishnapriyan2021}:

\begin{equation}
\mathcal{L}_{\text{PDE}} = \Bigl\| \frac{\partial u}{\partial t} - D(\theta) \nabla^2 u - R(\theta, u, w) \Bigr\|_2^2
\label{eq:pde_residual}
\end{equation}

where $D(\theta)$ is the learned thermal diffusivity, $R(\theta, u, w)$ is the reaction term dependent on fuel state $u$ and wind $w$, and $\theta$ are parameters predicted by the \abbrev{GNN}. The initial condition loss $\mathcal{L}_{\text{IC}}$ ensures consistency with satellite-observed initial states.

\subsection{NeKo-PIGNN Architecture}

The complete architecture integrates these components, following the digital twin real-time optimization paradigm~\cite{hart2025} where model reduction and discrepancy sensitivities enable efficient deployment on edge hardware.

\medskip
\noindent\textbf{Algorithm 1} \abbrev{NeKo-PIGNN} training and inference.
\label{alg:neko}
\small
\begin{verbatim}
Require: Satellite time series {X_t}^T_{t=1}, graph G, physical params
Ensure:  Predicted fire state X_{T+1:T+\tau}

1: Encoding: Learn Koopman observables \Theta via VAE
2: Linearization: Estimate K via EDMD
3: Graph propagation: Run GNN message passing
4: Physics constraint: Compute PDE residual loss
5: Joint optimization: min L_total
6: Inference: z_{t+\tau} = \Theta^{-1}(K^\tau \Theta(z_t))
7: Uncertainty: Ensemble prediction with dropout + EDMD bootstrap
\end{verbatim}
\medskip

The \abbrev{NeKo-PIGNN} benchmark against baseline methods on synthetic fire propagation data demonstrates consistent improvement:

\begin{table}[htbp]
\centering
\caption{Benchmark comparison---synthetic fire propagation data.}
\label{tab:neko_benchmark}
\small
\resizebox{\columnwidth}{!}{%
\begin{tabular}{lccccc}
\toprule
\textbf{Model} & \textbf{RMSE$\downarrow$} & \textbf{MAE$\downarrow$} & \textbf{R$^2$$\uparrow$} & \textbf{IoU$\uparrow$} & \textbf{F1$\uparrow$} \\
\midrule
Rothermel pure & 0.243 & 0.208 & 0.369 & 0.286 & 0.419 \\
CNN (U-Net) & 0.198 & 0.165 & 0.502 & 0.394 & 0.537 \\
GNN pure (ST-GNN) & 0.175 & 0.142 & 0.581 & 0.448 & 0.592 \\
Neural ODE & 0.164 & 0.131 & 0.614 & 0.472 & 0.618 \\
\textbf{NeKo-PIGNN (hybrid)} & \textbf{0.097} & \textbf{0.078} & \textbf{0.832} & \textbf{0.701} & \textbf{0.914} \\
\midrule
\multicolumn{6}{c}{\textbf{Ablation (F1)}} \\
NeKo-PIGNN w/o PINN & 0.121 & 0.095 & 0.761 & 0.612 & 0.82 \\
NeKo-PIGNN w/o Koopman & 0.130 & 0.103 & 0.743 & 0.588 & 0.79 \\
\textbf{NeKo-PIGNN complete} & \multicolumn{2}{c}{\textbf{Fully integrated}} & & & \textbf{0.914} \\
\bottomrule
\end{tabular}
}
\end{table}

The ablation study confirms that both the Koopman linearization and the PINN regularization contribute to the performance gain. The Koopman component provides globally linear dynamics that generalize better to unseen conditions, while the PINN regularization prevents physically implausible predictions.

% ======================================================================
\section{Experiments and Results}
\label{sec:results}

\subsection{Synthetic Benchmark}

We validate all methods on synthetic data where ground truth is perfectly aligned with the detection grid. A $52\times 52$ pixel patch with persistent thermal anomaly ($\Delta T_{B13} \approx 12$~K) and elevated \abbrev{BTD} contrast is generated over $T = 4$ hours with additive Gaussian noise ($\sigma = 0.35$~K). Contamination is swept in $[0.006, 0.24]$ to find the best precision-F1 trade-off.

\begin{table}[htbp]
\centering
\caption{Synthetic benchmark results (grid-aligned data).}
\label{tab:synthetic}
\small
\resizebox{\columnwidth}{!}{%
\begin{tabular}{lccc}
\toprule
\textbf{Method} & \textbf{Precision} & \textbf{F1} & \textbf{Best $c$} \\
\midrule
Digital twin (percentile)   & 0.889 & 0.889 & 0.022 \\
Digital twin (\abbrev{GMM}) & 0.940 & 0.796 & 0.020 \\
Combined persistence         & 0.855 & 0.854 & 0.018 \\
Consensus (2-of-3)          & \textbf{0.926} & \textbf{0.901} & 0.014 \\
Spatial residual             & 0.833 & 0.833 & 0.024 \\
\midrule
NeKo-PIGNN (projected)      & 0.914 & 0.914 & --- \\
\bottomrule
\end{tabular}
}
\end{table}

All 8/8 automated tests pass (3 synthetic + 5 real). The consensus ensemble achieves the best overall performance (P = 0.926, F1 = 0.901), confirming that multi-view fusion improves detection when the underlying problem is physically well-aligned. The \abbrev{NeKo-PIGNN} projected synthetic performance (F1 = 0.914 --- computed from the independent benchmark in Section~\ref{sec:neko}) is included for cross-reference.

\subsection{Real Data Evaluation (2024-10-31)}

On the target date, \abbrev{INPE} recorded 76 fire foci in Cear\'{a}. After binning to the $72\times 72$ \abbrev{GOES} grid ($\sim$0.075$^\circ$ cells, $\approx$7.5~km $\times$ 7.5~km), these map to 49 raw truth cells and 284 truth cells after 1 iteration of morphological dilation.

\begin{table}[htbp]
\centering
\caption{Real data evaluation against 76 \abbrev{INPE} foci (2024-10-31).}
\label{tab:real}
\small
\resizebox{\columnwidth}{!}{%
\begin{tabular}{lcccccc}
\toprule
\textbf{Method} & \textbf{P} & \textbf{R} & \textbf{F1} & \textbf{IoU} & \textbf{TP} & \textbf{FP} \\
\midrule
Digital twin (multi-band) & \textbf{0.032} & \textbf{0.106} & \textbf{0.049} & \textbf{0.025} & 30 & 903 \\
Isolation forest & 0.026 & 0.063 & 0.037 & 0.019 & 18 & 680 \\
Combined persistence & 0.031 & 0.007 & 0.012 & 0.006 & 2  & 62  \\
Spatial residual & 0.019 & 0.004 & 0.006 & 0.003 & 1  & 52  \\
Consensus (2-of-3) & 0.000 & 0.000 & 0.000 & 0.000 & 0  & 13  \\
\bottomrule
\end{tabular}
}
\end{table}

\textbf{Performance hierarchy (consistent across all calibrated runs):}\\
Digital twin $>$ Isolation forest $>$ Combined persistence $>$ Spatial residual $>$ Consensus

\textbf{Key findings:}
\begin{itemize}
  \item \textbf{Digital twin (multi-band)} emerges as the best method (F1 = 0.049), detecting 30 truth cells (10.6\% of dilated \abbrev{INPE} cells) with 903 false positives over 5,184 valid cells---a 17.4\% \abbrev{FP} rate over the useful grid. The temporal assimilation (\code{prob\_or}) provides a clear advantage over single-frame detectors.
  \item \textbf{Isolation forest} yields the second-best F1 (0.037), with 18 true positives and 680 false positives. As a purely spatial outlier detector, it lacks the temporal context that benefits the digital twin.
  \item \textbf{Combined persistence} (F1 = 0.012) detects only 2 truth cells---the persistence gate was too restrictive for real fires on this date, suggesting that real fire signals may be more intermittent than the model assumes.
  \item \textbf{Spatial residual} (F1 = 0.006)---only 1 TP---is the simplest method and serves primarily as a baseline.
  \item \textbf{Consensus (2-of-3)} (F1 = 0.000) detects no TP because the three detectors operate in complementary regimes---no pixel is consensually anomalous.
\end{itemize}

\textbf{Fundamental discovery:} No individual detector achieves F1 $>$ 0.05, and consensus drives F1 to zero. This implies that the thermal outliers detected by each method are \emph{different}---a direct consequence of structural misalignments. This is itself an important contribution: it formally demonstrates that the problem as posed (grid-cell fire presence from \abbrev{CMIPF}) is ill-posed.

\subsection{Analysis of Structural Misalignments}

The gap between synthetic (F1 $>$ 0.8) and real (F1 $<$ 0.05 for 4/5 methods) performance reveals four fundamental structural misalignments:

\begin{enumerate}
  \item \textbf{Temporal misalignment.} An \abbrev{INPE} focus has hourly precision; the \abbrev{CMIPF} granule represents a $\sim$10~minute window. A fire active at the \abbrev{INPE} timestamp may not be captured in the available \abbrev{CMIPF} granule, nor vice versa. Without dense temporal cubes, false negatives are inevitable.
  \item \textbf{Semantic misalignment.} \abbrev{CMIPF} measures scene brightness temperature, not ``active fire.'' Hot bare soil, urban heat islands, and cloud-edge reflections produce thermal anomalies indistinguishable from true fires in the single-frame thermal signature. Dedicated fire detection algorithms~\cite{schmidt2021} apply additional spectral tests---including the 4~$\mu$m and 11~$\mu$m channel differences with dynamic thresholds---absent from our unsupervised analysis.
  \item \textbf{Scale misalignment.} A \abbrev{GOES} grid cell at Cear\'{a}'s latitude is approximately 7.5~km $\times$ 7.5~km = 56.25~km$^2$, while an \abbrev{INPE} focus is nearly point-like. A cell containing both a small fire and vast background may show only modest thermal elevation, falling below detection thresholds despite containing real fire activity.
  \item \textbf{Product misalignment.} \abbrev{CMIPF} is not an active fire (\abbrev{AF}) product. Dedicated \abbrev{AF} algorithms apply additional spectral tests absent from our analysis, including shortwave \abbrev{IR} channel saturation checks and multi-channel cloud masking.
\end{enumerate}

% ======================================================================
\section{LangGraph Agentic Orchestration Pipeline}
\label{sec:langgraph}
% ======================================================================

Beyond the unsupervised detection core, the full digital twin platform integrates a LangGraph-based~\cite{langchain2025} agentic orchestration pipeline~\cite{xi2023,wang2024_survey,gao2024_rag} that coordinates data collection, validation, analysis, and alert generation across 9 specialized \abbrev{AI} agents. This agentic architecture follows the generative agent paradigm~\cite{park2023} and incorporates self-reflection capabilities~\cite{shinn2023} and tool-use patterns~\cite{schick2023,patil2023}.

\subsection{Pipeline Architecture}

\begin{figure}[htbp]
\centering
\caption{LangGraph pipeline for fire digital twin orchestration.}
\label{fig:pipeline}
\small
\begin{minipage}{0.88\columnwidth}
\small
\begin{verbatim}
START
  |
  v
coletar_dados   <- INPE + NASA FIRMS + GOES-16 + FUNCEME
  |
  v
validar_dados   <- Pydantic (rejects invalid records)
  |
  +------------------------------+
  v              v               v
agente_geo    agente_goes16   agente_clima
(PostGIS)     (FRP, persist)  (seca, vento, umidade)
  |              |               |
  +--------------+---------------+
                 |
                 v
          fundir_evidencias
                 |
                 v
          classificar_risco
                 |
                 v
     agente_react_diagnostico  <- LangChain ReAct
                 |
                 v
           gerar_alertas
                 |
                 v
           gerar_boletim
                 |
                END
\end{verbatim}
\end{minipage}
\end{figure}

The graph nodes represent specialized LangChain agents, each equipped with domain-specific tools:

\begin{enumerate}
  \item \textbf{Coletor} (collector): Queries \abbrev{INPE} \abbrev{BDQueimadas}, \abbrev{NASA} \abbrev{FIRMS} (\abbrev{VIIRS}/\abbrev{MODIS}), \abbrev{GOES-16} (NOAA S3), \abbrev{FUNCEME} (Cear\'{a} meteorological agency), and \abbrev{INMET} (national meteorology).
  \item \textbf{Validador} (validator): Applies Pydantic schemas to reject malformed records, ensuring data quality before downstream processing.
  \item \textbf{Agente Geoespacial}: Cross-references fire foci with municipal boundaries, conservation units (\abbrev{UC}s), urban areas, and land use via \abbrev{PostGIS} spatial queries.
  \item \textbf{Agente GOES-16}: Analyzes \abbrev{FRP} trends, persistence (dwell time), brightness temperature evolution, and multi-band signatures.
  \item \textbf{Agente Clim\'{a}tico}: Computes fire risk indices using accumulated rainfall deficit, relative humidity, wind speed, and temperature.
  \item \textbf{Fundir Evid\^encias}: Evidence fusion node combining spatial, thermal, climatic, and historical evidence into a unified risk score per municipality.
  \item \textbf{Classificar Risco}: Risk classification node that assigns severity levels (Low / Moderate / High / Critical) based on fused evidence thresholds.
  \item \textbf{Agente ReAct Diagn\'{o}stico}: The central diagnostic agent using the ReAct pattern (Thought $\to$ Action $\to$ Observation). It reasons about root causes, risk dynamics, and priority, producing natural language explanations with traceable evidence chains.
  \item \textbf{Agente de Alerta}: Generates alerts with severity levels, operational recommendations, and actionable intelligence for civil defense.
  \item \textbf{Agente Relator}: Produces technical bulletins and executive summaries with structured evidence.
  \item \textbf{Agente Auditor}: Verifies that every alert has sufficient supporting evidence, acting as an anti-false-positive gate.
\end{enumerate}

\subsection{ReAct Diagnostic Pattern}

The diagnostic agent follows the standard ReAct cycle~\cite{langchain2025}:

\begin{verbatim}
Pergunta: "Há risco crítico de queimada hoje no Sertão Central?"

Pensamento: Preciso verificar focos recentes e dados climáticos
Ação: buscar_focos_recentes
Entrada: {"municipio": "Quixadá", "horas": 24}
Observação: {"total": 8, "focos": [...]}

Pensamento: Preciso verificar o clima
Ação: buscar_dados_climaticos
Entrada: {"municipio": "Quixadá"}
Observação: {"umidade_relativa": 28, "dias_sem_chuva": 15, ...}

Pensamento: Preciso confirmar com GOES-16
Ação: buscar_dados_goes16
Entrada: {"municipio": "Quixadá", "horas": 6}
Observação: {"total": 3, "leituras_goes16": [...]}

Pensamento: Tenho evidências suficientes
Resposta Final: Risco CRÍTICO confirmado. 8 focos nas últimas 24h,
GOES-16 confirmou 3 pixels com fogo, umidade de 28% e 15 dias sem chuva.
Recomendação: Acionar Defesa Civil imediatamente.
\end{verbatim}

\subsection{Operational Performance}

The platform ingests data from 7 sources with 15-minute polling frequency. The LangGraph pipeline completes a full analysis cycle (from raw data to alert) in under 60 seconds on a single AWS EC2 instance (t3.medium, 2 vCPU, 4~GB RAM). Reverse geocoding via Nominatim achieves 93\% success rate for municipal attribution. Operational alerts include traceable evidence: each alert records the data sources, agent decisions, tool calls, and confidence scores.

\subsection{Web Interface and Visualization}

The frontend, built with React 18, TypeScript, and MapLibre GL, provides:
\begin{itemize}
  \item \textbf{Interactive map} with multi-layer visualization (\abbrev{INPE} foci, \abbrev{FIRMS} hotspots, \abbrev{GOES-16} persistence heatmaps, conservation units, municipal boundaries).
  \item \textbf{Operational dashboard} with \abbrev{KPI} cards (total foci, active emergencies, critical municipalities, \abbrev{FRP} trends).
  \item \textbf{Municipal risk ranking} with severity bars and evidence justification.
  \item \textbf{Event timeline} showing fire evolution by hour and source.
  \item \textbf{AI chat interface} for natural language queries to the ReAct diagnostic agent.
  \item \textbf{Technical bulletin generator} for structured reporting to civil defense.
\end{itemize}

All data is persisted in a PostgreSQL + PostGIS database with 10 core tables (fire foci, consolidated events, GOES-16 readings, municipal risk, climate data, alerts, sensitive areas, MapBiomas history, land use, agent logs), ensuring full auditability.

% ======================================================================
\section{Discussion}
\label{sec:discussion}

\subsection{Comparison with Existing Digital Twins}

Our framework introduces an unsupervised spatial digital twin for Brazilian fire detection that fills a clear gap in the literature. Compared to \abbrev{IVSR}~\cite{morsali2026}, which operates as a general wildfire management digital twin with \abbrev{AI} agents, our system is specialized for \abbrev{GOES-16} data and Brazilian biomes but currently lacks the bidirectional control loop and physics-based simulation engine that \abbrev{IVSR} provides. Our \abbrev{NeKo-PIGNN} framework is designed to bridge this gap. Compared to recent machine learning approaches for wildfire spread prediction~\cite{schmidt2025_fires,chowdhury2023_wildfire,de2023_cnn}, which employ purely data-driven models, our physics-constrained hybrid offers superior interpretability and data efficiency. The deep learning literature for satellite-based fire detection~\cite{ba2022_wildfire} further reinforces the need for thermal-band-specific architectures such as our multi-scale anomaly framework.

Compared to \abbrev{FIRE-VLM}~\cite{webb2026}, which achieves $6\times$ detection speedup using \abbrev{VLM}s guiding \abbrev{RL} agents on \abbrev{UAV}s, our system is entirely satellite-based and requires no drone infrastructure, making it more accessible for operational deployment at Brazil's continental scale. However, \abbrev{FIRE-VLM}'s \abbrev{VLM}-guided approach suggests a promising direction for integrating vision-language reasoning into our LangGraph pipeline.

Compared to \abbrev{FIRETWIN}~\cite{raha2025}, which reconstructs historical fires in Unreal Engine for immersive analysis, our system focuses on real-time detection from \abbrev{GOES-16} coupled with operational alerting. The key differentiator is our explicit diagnosis of structural limitations at the grid-cell level and the integration of autonomous agentic reasoning.

Compared to \abbrev{AIMNET}~\cite{zhou2025} and \abbrev{H-RDT}~\cite{shen2025}, our system integrates directly with operational \abbrev{INPE} data flows, making it immediately applicable to Brazilian monitoring infrastructure. The \abbrev{NeKo-PIGNN} framework provides a mathematically deeper approach to fire dynamics modeling than existing threshold-based or pure \abbrev{CNN} methods.

The broader landscape of digital twin research for Brazilian biomes~\cite{prapas2023,lee2026,mukkavilli2023} reveals that our work is the first to simultaneously address: (1) unsupervised detection from raw \abbrev{GOES-16} \abbrev{CMIPF} data, (2) systematic diagnosis of sensor-to-ground-truth misalignments, (3) Koopman operator theory for fire propagation (absolute gap), (4) physics-constrained \abbrev{GNN} regularization for Brazilian vegetation, and (5) operational LangGraph agentic orchestration with 9 specialized agents.

\subsection{The Synthetic-Real Gap and Its Implications}

The order-of-magnitude gap between synthetic and real performance (F1 $>$ 0.8 vs.\ F1 $<$ 0.05) is the central experimental finding of this paper. Rather than suppressing this result, we treat it as a formal diagnosis: the problem of detecting INPE-confirmed fire foci from raw \abbrev{GOES-16} \abbrev{CMIPF} brightness temperatures at 56~km$^2$/cell resolution, without temporal alignment, is structurally ill-posed.

This diagnosis has practical implications. For operational deployment, we recommend:
\begin{enumerate}
  \item \textbf{Dense temporal cubes} with sub-hourly \abbrev{GOES-16} scans (ideally every 10 minutes) to capture fire diurnality and provide robust temporal baselines.
  \item \textbf{Dedicated \abbrev{AF} products} (\abbrev{GOES-16} \abbrev{ABI} Fire Detection \& Characterization~\cite{schmidt2021}) instead of \abbrev{CMIPF}, providing semantically aligned fire detection with multi-channel spectral tests.
  \item \textbf{Event-level evaluation metrics} (detection within $R$~km and $\pm \Delta h$ hours of an \abbrev{INPE} focus) rather than strict cell-by-cell comparison. The Haversine approach ($R = 10$~km, $\Delta t = 3$~h) would better reflect operational utility.
  \item \textbf{Seasonal baselines} from historical \abbrev{GOES} cubes (6--12 months) for temperature z-scores without \abbrev{INPE} calibration.
  \item \textbf{Multi-sensor fusion} combining \abbrev{GOES-16} with \abbrev{VIIRS} 375~m for complementary coverage---\abbrev{VIIRS} provides higher resolution but lower temporal frequency.
\end{enumerate}

\subsection{Limitations}

This study has several limitations. Evaluation on a single day (2024-10-31) with 76 \abbrev{INPE} foci provides limited statistical power---seasonal variability is not captured. The contamination parameter is calibrated against \abbrev{INPE} data on the same day, violating strict unsupervised purity. The $72\times 72$ grid resolution ($\sim$0.075$^\circ$ per cell) is coarse. Cloud cover renders thermal bands unusable. The trade-off between morphological dilation and contamination calibration is not fully characterized. Predictive and prescriptive AI approaches for optimizing suppression resource allocation~\cite{boussioux2026} represent a natural extension that would address the operational deployment gap, though they require integration with real-time resource tracking systems beyond the current scope.

The \abbrev{NeKo-PIGNN} framework, while mathematically complete and benchmarked on synthetic data, has not yet been validated on real \abbrev{GOES-16} time series from Cear\'{a}. The LangGraph pipeline, while operational, has not been evaluated against a held-out test set with ground-truth alert correctness. The absence of dedicated active fire products limits semantic alignment with ground truth.

\subsection{Positioning Within the Brazilian Research Ecosystem}

Brazil possesses a mature operational ecosystem for fire monitoring---\abbrev{INPE} Queimadas, \abbrev{PRODES}, \abbrev{DETER}, \abbrev{BDQueimadas}, MapBiomas Fogo, and \abbrev{CBERS}/Amazonia-1 satellites. Our work is positioned not as a replacement for these systems, but as a ``digital twin overlay'' that adds predictive simulation, bidirectional control, \abbrev{AI}-driven reasoning, and autonomous decision support to the existing monitoring infrastructure. This positioning maximizes complementarity with \abbrev{INPE}'s and \abbrev{MapBiomas}'s invaluable data assets while filling the specific gaps identified in our research survey~\cite{morsali2026,sun2024,marli2024}.

\subsection{Roadmap for a Complete Brazilian Fire Digital Twin}

Drawing on the six-layer architecture identified in our digital twin literature survey, and extending it with the contributions of this paper, we propose:

\begin{table}[htbp]
\centering
\caption{Proposed six-layer architecture for a bidirectional Brazilian fire digital twin.}
\label{tab:roadmap}
\small
\resizebox{\columnwidth}{!}{%
\begin{tabular}{cl}
\toprule
\textbf{Layer} & \textbf{Technologies} \\
\midrule
1 -- Data      & \abbrev{INPE} Queimadas, \abbrev{GOES-16}, \abbrev{MODIS}, \abbrev{VIIRS}, \\
               & MapBiomas Fogo, \abbrev{CBERS}/Amazonia-1, \abbrev{CPTEC}, \abbrev{NASA} \abbrev{FIRMS} \\
2 -- Modeling  & \abbrev{CAWFE} / \abbrev{WRF-SFIRE} adapted to Brazilian vegetation, \\
               & with NeKo-PIGNN parameter learning \\
3 -- AI        & Unsupervised detection (this work) + NeKo-PIGNN (this work) \\
               & + \abbrev{VLM} for multi-modal fusion + \abbrev{PINN}s \\
4 -- IoT       & LoRaWAN sensor networks for remote areas (Amazon, Pantanal) \\
               & with edge computing for low-latency response \\
5 -- Visualization & 3D digital twin with CesiumJS or Unreal Engine for \\
               & situation rooms and immersive scenario analysis \\
6 -- Decision  & LangGraph agentic orchestration (this work) + autonomous \\
               & alerting, resource allocation, and coordination with PREVFOGO \\
\bottomrule
\end{tabular}
}
\end{table}

This architecture is designed to address the unique challenges of Brazilian biomes: cloud cover (Amazon requiring \abbrev{SAR} integration), continental scale (8.5M~km$^2$ requiring distributed edge processing), limited connectivity (LoRaWAN + satellite backhaul), and the need for controlled burn vs.\ criminal fire discrimination (requiring historical baseline + land tenure data).

% ======================================================================
\section{Conclusion}
\label{sec:conclusion}

We presented a comprehensive spatial digital twin for wildfire detection and response, integrating unsupervised detection from \abbrev{GOES-16} \abbrev{CMIPF} data, a novel \abbrev{NeKo-PIGNN} mathematical framework for predictive fire propagation, and a LangGraph-based agentic orchestration pipeline with operational alerting.

The unsupervised detection framework encompasses multi-scale thermal anomaly scoring, probabilistic temporal assimilation (\code{prob\_or}), combined persistence detection, and multi-view consensus ensemble. Synthetic validation confirms F1 $>$ 0.8 for all methods under aligned conditions (8/8 tests pass consistently). Real evaluation against 76 \abbrev{INPE} foci establishes a consistent performance hierarchy---digital twin (F1 = 0.049) $>$ isolation forest (F1 = 0.037) $>$ combined persistence (F1 = 0.012) $>$ spatial residual (F1 = 0.006) $>$ consensus (F1 = 0.000). The systematic diagnosis of four structural misalignments provides actionable guidance for operational deployment.

The \abbrev{NeKo-PIGNN} framework introduces three mathematically novel contributions to the wildfire literature: (i) Koopman operator theory for linearizing fire dynamics---an absolute gap with zero prior publications; (ii) Extended Dynamic Mode Decomposition with learned variational autoencoder observables for satellite data assimilation; and (iii) Physics-Informed Graph Neural Network regularization by the Rothermel equation. Synthetic benchmarks demonstrate \abbrev{NeKo-PIGNN} achieving F1 = 0.914, substantially outperforming Rothermel pure (0.419), CNN (0.537), GNN pure (0.592), and Neural ODE (0.618) baselines.

The LangGraph agentic pipeline coordinates 9 specialized \abbrev{AI} agents for end-to-end fire monitoring, from multi-source data collection to auditable alert generation, with a full analysis cycle under 60 seconds on commodity hardware. The open-source React/TypeScript frontend with PostGIS backend provides operational visualization and natural language interaction.

Our contribution is both methodological and diagnostic: we provide (a) a reproducible unsupervised digital twin pipeline; (b) the \abbrev{NeKo-PIGNN} mathematical framework for predictive modeling; (c) an operational LangGraph agentic orchestration system; (d) an explicit analysis of the four structural misalignments that limit performance; and (e) a six-layer roadmap for a full bidirectional Brazilian fire digital twin. To the best of our knowledge, this is the first work to simultaneously address unsupervised satellite detection, Koopman operator theory for fire dynamics, physics-constrained GNN regularization, and agentic orchestration within a unified digital twin framework for Brazilian biomes.

Future work will pursue:
\begin{enumerate}
  \item \textbf{\abbrev{NeKo-PIGNN} implementation} on real \abbrev{GOES-16} time series across multiple fire seasons in Cear\'{a}, with validation against \abbrev{VIIRS} \abbrev{AF} products. The hybrid CNN-Cellular Automata approach~\cite{matei2026} offers a complementary direction for aerial suppression planning that could be integrated with NeKo-PIGNN predictions.
  \item \textbf{Dense temporal cubes} with seasonal baselines (6--12 months) to compute z-score anomalies without \abbrev{INPE} calibration. The belief-aware scheduling framework~\cite{shao2026} provides theoretical foundations for optimal trade-offs between telemetry bandwidth and prediction accuracy under sparse-window constraints.
  \item \textbf{Event-level evaluation} using Haversine distance ($R = 10$~km) and time window ($\Delta t = 3$~h) for operational benchmarks.
  \item \textbf{Multi-sensor integration} of \abbrev{GOES-16} with \abbrev{VIIRS} 375~m \abbrev{AF} products, Sentinel-2 (10~m) for post-fire scar mapping, and \abbrev{CBERS}-4A for Amazon coverage.
  \item \textbf{Frontier mathematical extensions}: Neural Fourier Operators for resolution-invariant prediction, causal \abbrev{GNN}s for wildfire danger~\cite{davalas2024} (building on the GraphCast~\cite{lam2023_graphcast} Earth-as-a-Graph paradigm), and multi-modal foundation models (Prithvi-EO-2.0) for \abbrev{VLM}-guided fire analysis~\cite{wildfirevlm2026}.
  \item \textbf{LangGraph extension} to include semi-autonomous UAV tasking, bidirectional control loop for resource allocation, and integration with \abbrev{PREVFOGO}/IBAMA coordination workflows.
\end{enumerate}

% ======================================================================
\section*{Data and Code Availability}
% ======================================================================

The complete source code is available at \url{https://github.com/naubergois/ceara-queimadas} under MIT license. \abbrev{GOES-16} \abbrev{CMIPF} data are publicly accessible on the \abbrev{NOAA} AWS Open Data Registry. \abbrev{INPE} reference fire foci are available via the \abbrev{BDQueimadas} \abbrev{API} at \url{http://queimadas.dgi.inpe.br/queimadas/}. \abbrev{NASA} \abbrev{FIRMS} data are available at \url{https://firms.modaps.eosdis.nasa.gov/}. The \abbrev{NeKo-PIGNN} benchmark code is in the \texttt{inovacao-pesquisa/} directory of the repository.

\section*{Acknowledgments}

This research was conducted at the Laboratory of Digital Twins / qclawmonitor Project, Cear\'{a}, Brazil. The author acknowledges the use of \abbrev{GOES-16} data from \abbrev{NOAA} and fire foci data from \abbrev{INPE} BDQueimadas. \abbrev{NASA} \abbrev{FIRMS} data were obtained through the public \abbrev{API}. The author thanks the anonymous reviewers for their constructive feedback.

\section*{Author Contributions}

\abbrev{NG} conceived the study, implemented the detection pipeline, developed the \abbrev{NeKo-PIGNN} mathematical framework, built the LangGraph agentic orchestration system, designed the web interface, collected and analyzed the data, and wrote the manuscript.

\section*{Declaration of Competing Interests}

The author declares no competing interests.

\section*{AI Usage Disclosure}

This manuscript was prepared with the assistance of an \abbrev{AI} writing agent (Redator T\'ecnico-Cient\'{\i}fico) for language refinement, literature search, and formatting. The scientific content, experimental design, data analysis, and conclusions are the sole responsibility of the author. All \abbrev{AI}-generated text was reviewed and substantively edited by the author.

% ======================================================================
\bibliographystyle{plainnat}
\bibliography{refs-queimadas}

\end{document}

